\chapter{Compacidad en espacios métricos}\label{cap:nuevo05}\label{chap:13}\label{chap:14}\label{chap:15}

\section{De los controles locales a la compacidad}\label{lectura:13}
Una función continua controla su variación cerca de cada punto, pero el tamaño de esa vecindad puede cambiar de un punto a otro. La compacidad permite reducir una familia de controles locales a una cantidad finita y, después, tomar el mínimo de los radios o el máximo de las cotas. Esta es la razón de su presencia en tantos argumentos de análisis. Usaremos abiertos relativos, continuidad, sucesiones de Cauchy y cerradura. Partiremos de las cubiertas abiertas, pasaremos a la extracción de subsucesiones y obtendremos después el criterio de completitud y acotación total. La compacidad relativa y la compacidad local mostrarán cómo usar estas ideas dentro de un espacio ambiente y alrededor de cada punto. Este capítulo integra las lecturas 13--15 de \cite{valencia}.

\section{Cubiertas y conjuntos compactos}
\begin{definicion}
Una \emph{cubierta abierta}\index{cubierta abierta} de $X$ es una familia $\mathcal U$ de abiertos de $X$ cuya unión es $X$. Una \emph{subcubierta} es una subfamilia que todavía cubre $X$. El espacio $X$ es \emph{compacto}\index{compacidad} si toda cubierta abierta tiene una subcubierta finita. Un subconjunto $K\subset X$ es compacto cuando lo es con su topología relativa. El conjunto vacío es compacto: la subfamilia vacía basta para cubrirlo.
\end{definicion}
Para un radio fijo $r>0$, las bolas $B(x,r)$, $x\in X$, cubren $X$. Si $D$ es denso, también las bolas con centros en $D$: dado $x\in X$, existe $a\in D\cap B(x,r)$ y la simetría de la distancia da $x\in B(a,r)$. En $\R$, los intervalos $(n,\infty)$, $n\in\Z$, cubren la recta pero ninguna cantidad finita la cubre. En un espacio discreto la familia pertinente es $\{\{x\}:x\in X\}$; sus elementos son conjuntos, mientras que los elementos de $X$ son puntos. Esta cubierta no admite quitar un solo miembro.

\begin{proposicion}\label{ch13:criterios}
Sean $\mathcal B$ una base de la topología de $X$ y $K\subset X$.
\begin{enumerate}[label=(\roman*)]
\item $X$ es compacto si y solo si toda cubierta de $X$ formada por miembros de $\mathcal B$ admite una subcubierta finita.
\item $K$ es compacto si y solo si toda familia de abiertos de $X$ que cubre $K$ admite una subfamilia finita que cubre $K$.
\end{enumerate}
\end{proposicion}
\begin{proof}
Para (i), una dirección es parte de la definición. Para la otra, dada una cubierta abierta $\mathcal U$, tomamos todos los $B\in\mathcal B$ que estén contenidos en algún $U\in\mathcal U$. Cubren $X$: si $x\in U$, la propiedad de base proporciona $x\in B\subset U$. Una subcubierta $B_1,\ldots,B_m$ permite elegir $U_i$ con $B_i\subset U_i$, y estos $U_i$ cubren $X$.

Para (ii), si $K$ es compacto, los abiertos relativos $K\cap U$ forman una cubierta de $K$ y por tanto basta una cantidad finita. Recíprocamente, toda cubierta relativa tiene la forma $\{K\cap U_\alpha\}$ con $U_\alpha$ abierto en $X$. Aplicar la hipótesis a $\{U_\alpha\}$ da la subcubierta relativa requerida.
\end{proof}

\begin{ejemplo}[Espacios finitos e intervalos cerrados]
Todo espacio finito es compacto: se elige un miembro de la cubierta para cada punto. Un espacio discreto infinito no lo es, por su cubierta de singletons.

El intervalo $[a,b]$ es compacto. Presentamos la prueba, en lugar de dar por conocida esta propiedad del cálculo. Sea $\mathcal U$ una cubierta y definamos
\[
E=\{t\in[a,b]:[a,t]\text{ se cubre con finitos miembros de }\mathcal U\}.
\]
El conjunto contiene $a$ y tiene supremo $c\in[a,b]$. Si $c<b$, un abierto $U\ni c$ contiene $(c-\delta,c+\delta)\cap[a,b]$ para algún $\delta>0$. Por la definición de supremo hay $t\in E$ con $t>c-\delta/2$; si $c=a$ podemos tomar $t=a$. Una cubierta finita de $[a,t]$, junto con $U$, cubre $[a,\min(b,c+\delta/2)]$, contradiciendo que $c$ sea cota superior. Así $c=b$. Tomando ahora un abierto que contenga $b$ y un $t\in E$ suficientemente próximo a $b$, obtenemos una cubierta finita de todo $[a,b]$.
\end{ejemplo}

\section{Imágenes, cerrados y continuidad uniforme}
\begin{teorema}\label{ch13:imagen}
La imagen continua de un espacio compacto es compacta.
\end{teorema}
\begin{proof}
Sea $f:X\to Y$ continua y $\mathcal U$ una familia de abiertos de $Y$ que cubre $f(X)$. Las preimágenes $f^{-1}(U)$ son abiertas y cubren $X$. Elegimos $U_1,\ldots,U_m$ cuyas preimágenes cubran $X$. Si $y=f(x)$, entonces $x\in f^{-1}(U_i)$ para algún $i$, de donde $y\in U_i$. La proposición \ref{ch13:criterios} termina la prueba.
\end{proof}

\begin{proposicion}\label{ch13:cerrados}
En un espacio métrico, todo compacto es cerrado y acotado. Todo subconjunto cerrado de un compacto es compacto.
\end{proposicion}
\begin{proof}
Sea $K$ compacto no vacío y $x\notin K$. Para $y\in K$ pongamos $r_y=d(x,y)/3>0$. Elegimos $y_1,\ldots,y_m$ tales que $K\subset\bigcup_iB(y_i,r_i)$ y ponemos $r=\min_i r_i>0$. Si $z\in B(x,r)$ y $y\in K$, elegimos $i$ con $d(y,y_i)<r_i$. La desigualdad triangular, escrita en sentido inverso, da
\[
d(z,y)\ge d(x,y_i)-d(x,z)-d(y,y_i)>3r_i-r-r_i\ge r.
\]
Así $z\notin K$ y $B(x,r)\subset X\setminus K$. Esto prueba que $K$ es cerrado. Para la acotación, fijemos $x_0\in X$; las bolas $B(x_0,n)$, $n\in\N$, cubren $K$. Una subcubierta finita queda contenida en la de mayor radio.

Si $F$ es cerrado en el compacto $K$, una cubierta relativa de $F$ por abiertos de $K$, junto con $K\setminus F$, cubre $K$. Extraemos una subcubierta finita y eliminamos $K\setminus F$ si aparece. Los miembros restantes cubren $F$.
\end{proof}

\begin{teorema}[Heine--Cantor]\label{ch13:heinecantor}
Toda función continua $f:X\to Y$ entre espacios métricos, con $X$ compacto, es uniformemente continua.
\end{teorema}
\begin{proof}
Fijemos $\eps>0$. Para cada $x\in X$ elegimos $r_x>0$ tal que $d(x,z)<r_x$ implique $d_Y(f(x),f(z))<\eps/2$. Extraemos una subcubierta finita $B(x_i,r_i/2)$ y ponemos $\delta=\frac12\min_i r_i$. Dados $u,v$ con $d(u,v)<\delta$, elegimos $i$ con $u\in B(x_i,r_i/2)$. Entonces
\[
d(v,x_i)\le d(v,u)+d(u,x_i)<\delta+r_i/2\le r_i.
\]
Las imágenes de $u$ y $v$ están a distancia menor que $\eps/2$ de $f(x_i)$, por lo que $d_Y(f(u),f(v))<\eps$. El radio $\delta$ se eligió antes que $u,v$: esta independencia es exactamente la continuidad uniforme.
\end{proof}

\begin{corolario}[Máximo y mínimo]\label{ch13:extremos}
Si $K$ es compacto no vacío y $f:K\to\R$ es continua, existen $x_-,x_+\in K$ tales que $f(x_-)\le f(x)\le f(x_+)$ para todo $x\in K$.
\end{corolario}
\begin{proof}
La imagen es compacta, luego cerrada y acotada. Si $M=\sup f(K)$, para cada $n$ existe $y_n\in f(K)$ con $M-1/n<y_n\le M$. Por cerradura, $M\in f(K)$. Aplicamos el mismo argumento a $-f$ para el mínimo.
\end{proof}

\section{Ejercicios propuestos: cubiertas e imágenes compactas}
\begin{ejercicio}[Lectura 13, ejercicio 1]\label{ej:lectura-13-1}
Una familia no vacía $\mathcal F$ tiene la propiedad de intersección finita cuando toda intersección de una cantidad finita de sus miembros es no vacía. Probar que $X$ es compacto si y solo si toda familia de cerrados con esta propiedad tiene intersección no vacía.
\end{ejercicio}


\begin{ejercicio}[Lectura 13, ejercicio 2]\label{ej:lectura-13-2}
Probar que una unión finita de subconjuntos compactos es compacta.
\end{ejercicio}


\begin{ejercicio}[Lectura 13, ejercicio 3]\label{ej:lectura-13-3}
En un compacto $X$, una sucesión decreciente de cerrados no vacíos $F_n$ tiene intersección no vacía. Deducir que $X$ es completo.
\end{ejercicio}


\begin{ejercicio}[Lectura 13, ejercicio 4]\label{ej:lectura-13-4}
Probar que la bola unitaria cerrada de $\ell^\infty$ es cerrada y acotada, pero no compacta.
\end{ejercicio}


\begin{ejercicio}[Lectura 13, ejercicio 5]\label{ej:lectura-13-5}
Si $X$ es compacto y $f:X\to Y$ es continua e inyectiva, probar que $f^{-1}:f(X)\to X$ es continua.
\end{ejercicio}


\begin{ejercicio}[Lectura 13, ejercicio 6]\label{ej:lectura-13-6}
Sean $X$ compacto, $f:X\to Y$ continua y $A\subset X$ tal que $f(X\setminus A)\subset Y\setminus f(A)$. Probar: (i) $f|_A:A\to f(A)$ es cerrada; (ii) si además es inyectiva, es un homeomorfismo.
\end{ejercicio}


\begin{ejercicio}[Lectura 13, ejercicio 7]\label{ej:lectura-13-7}
Si $X,Y$ son compactos, probar que $X\times Y$, con la métrica máximo, es compacto.
\end{ejercicio}


\begin{ejercicio}[Lectura 13, ejercicio 8]\label{ej:lectura-13-8}
Probar que toda cubierta abierta $\mathcal U$ de un compacto métrico admite un número de Lebesgue: existe $r>0$ tal que todo subconjunto de diámetro menor que $r$ está contenido en algún miembro de $\mathcal U$.
\end{ejercicio}


La compacidad no significa simplemente que un conjunto sea pequeño en diámetro. Significa que cualquier descripción local mediante abiertos puede reducirse a finitos datos. La extracción de esa subfamilia finita es el paso que produce controles uniformes, extremos y continuidad de inversas.

\section{Compacidad secuencial y puntos de acumulación}\label{lectura:14}
Las cubiertas abiertas expresan bien el significado de compacidad, pero en problemas concretos solemos disponer de sucesiones: aproximaciones, minimizantes o funciones que queremos hacer converger. En espacios métricos ambas perspectivas son equivalentes. La palabra \emph{métrico} es importante: fuera de esta categoría no se pueden intercambiar sin más compacidad y compacidad secuencial. Los resultados sobre cubiertas que acabamos de demostrar permiten ahora construir límites de subsucesiones. Usaremos también puntos de acumulación y bases numerables.

\begin{definicion}
Un espacio es \emph{secuencialmente compacto}\index{compacidad!secuencial} si toda sucesión tiene una subsucesión convergente a un punto del espacio. Es \emph{compacto por puntos de acumulación}\index{compacidad!por puntos de acumulacion@por puntos de acumulación} si todo subconjunto infinito tiene un punto de acumulación en el espacio.
\end{definicion}

\begin{lema}\label{ch14:acumulacion}
Si $a$ es punto de acumulación de $A$ en un espacio métrico, cada bola alrededor de $a$ contiene infinitos puntos de $A$ distintos de $a$.
\end{lema}
\begin{proof}
Si $B(a,r)\cap(A\setminus\{a\})=\{b_1,\ldots,b_m\}$ fuera finito, podríamos elegir $0<s<r$ menor que todas las distancias positivas $d(a,b_j)$. Entonces $B(a,s)$ no contendría ningún punto de $A$ distinto de $a$, contradiciendo la acumulación. Si la lista ya es vacía, la contradicción es directa.
\end{proof}

\begin{teorema}\label{ch14:equivalencias}
Para un espacio métrico $X$ son equivalentes:
\begin{enumerate}[label=(\roman*)]
\item $X$ es compacto;
\item $X$ es secuencialmente compacto;
\item todo subconjunto infinito de $X$ tiene un punto de acumulación en $X$.
\end{enumerate}
\end{teorema}
\begin{proof}
\emph{(i)$\Rightarrow$(ii).} Dada $(x_n)$, ponemos $F_n=\cl{\{x_k:k\ge n\}}$. Son cerrados no vacíos y decrecientes. El ejercicio \ref{ej:lectura-13-3} proporciona $x\in\bigcap_nF_n$. Elegimos $n_1$ con $d(x_{n_1},x)<1$. Una vez elegido $n_j$, la pertenencia $x\in F_{n_j+1}$ permite escoger $n_{j+1}>n_j$ con $d(x_{n_{j+1}},x)<1/(j+1)$. Esta es una subsucesión que converge a $x$. La barra de cerradura en $F_n$ es indispensable.

\emph{(ii)$\Rightarrow$(iii).} Si $A$ es infinito, elegimos puntos distintos $a_n\in A$. Una subsucesión converge a cierto $a\in X$. Toda bola alrededor de $a$ contiene términos de esa subsucesión distintos de $a$, porque los términos son distintos entre sí y a lo sumo uno puede ser igual a $a$. Así $a$ es punto de acumulación.

\emph{(iii)$\Rightarrow$(ii).} Dada $(x_n)$, si su recorrido es finito, algún valor aparece infinitas veces: una unión finita de conjuntos finitos de índices no podría ser $\N$. Se obtiene una subsucesión constante. Si el recorrido es infinito, tiene un punto de acumulación $x$. Elegido $n_j$, el lema \ref{ch14:acumulacion} permite hallar en $B(x,1/(j+1))$ un valor del recorrido distinto de $x_1,\ldots,x_{n_j}$. Algún índice que lo produzca debe ser mayor que $n_j$. Se construye así una subsucesión convergente.

\emph{(ii)$\Rightarrow$(i).} Primero demostraremos dos hechos. Para cada $r>0$, finitas bolas de radio $r$ cubren $X$. Si esto fallara, elegiríamos $x_1$ y, sucesivamente, $x_{n+1}\notin\bigcup_{j=1}^nB(x_j,r)$. Las distancias entre términos distintos serían al menos $r$, de modo que ninguna subsucesión sería de Cauchy y ninguna convergería.

Segundo, toda cubierta abierta $\mathcal U$ tiene un número de Lebesgue. En caso contrario, para cada $n$ habría un conjunto no vacío $A_n$ de diámetro menor que $1/n$ que no estuviera contenido en ningún miembro de $\mathcal U$. Elegimos $a_n\in A_n$ y una subsucesión $a_{n_j}\to a$. Tomamos $U\in\mathcal U$ y $r>0$ con $B(a,r)\subset U$. Para $j$ grande, $d(a_{n_j},a)<r/2$ y $1/n_j<r/2$. Todo $z\in A_{n_j}$ satisface $d(z,a)<r$, luego $A_{n_j}\subset U$, contradicción.

Sea $\lambda$ un número de Lebesgue. Cubrimos $X$ con finitas bolas de radio $\lambda/3$. Cada una tiene diámetro a lo sumo $2\lambda/3<\lambda$, por lo que está contenida en algún miembro de $\mathcal U$. Estos miembros forman una subcubierta finita. Esta prueba no usa el lema de Lebesgue para compactos para demostrar su propia hipótesis: acabamos de obtener el lema directamente de la compacidad secuencial.
\end{proof}

\begin{corolario}\label{ch14:separable}
Todo compacto métrico es separable y admite una base numerable.
\end{corolario}
\begin{proof}
Para cada $n$ elegimos un conjunto finito $D_n$ cuyos centros de bolas de radio $1/n$ cubran $X$. La unión $D=\bigcup_nD_n$ es a lo sumo numerable. Si $B(x,r)$ es una bola, elegimos $n$ con $1/n<r$ y un centro $a\in D_n$ con $d(x,a)<1/n$. Así $D$ es denso. Las bolas de centro en $D$ y radio racional positivo forman una base: si $B(x,r)$ está dado, elegimos $a\in D$ suficientemente próximo a $x$ y un racional $q$ con $d(a,x)<q<r-d(a,x)$. Entonces $x\in B(a,q)\subset B(x,r)$.
\end{proof}

\section{Ejercicios propuestos: sucesiones, productos y colas}
\begin{ejercicio}[Lectura 14, ejercicio 1]\label{ej:lectura-14-1}
Si cada $(X_n,d_n)$ es compacto, probar que $\prod_{n\ge1}X_n$ es compacto para la métrica
\[
D(x,y)=\sum_{n=1}^\infty 2^{-n}\frac{d_n(x_n,y_n)}{1+d_n(x_n,y_n)}.
\]
\end{ejercicio}


\begin{ejercicio}[Lectura 14, ejercicio 2]\label{ej:lectura-14-2}
Sea $X$ compacto y $U\subset X\times X$ abierto que contiene la diagonal $\Delta_X$. Probar que existe $r>0$ tal que $d(x,y)<r$ implica $(x,y)\in U$.
\end{ejercicio}


\begin{ejercicio}[Lectura 14, ejercicio 3]\label{ej:lectura-14-3}
Probar que todo espacio métrico no compacto contiene un subconjunto infinito numerable, cerrado y discreto.
\end{ejercicio}


\begin{ejercicio}[Lectura 14, ejercicio 4]\label{ej:lectura-14-4}
Probar que $X$ es compacto si y solo si toda función continua $X\to\R$ es acotada.
\end{ejercicio}


\begin{ejercicio}[Lectura 14, ejercicio 5]\label{ej:lectura-14-5}
Sea $1\le p<\infty$. Probar que un subconjunto cerrado y acotado $A\subset\ell^p$ es compacto si y solo si
\[
\forall\eps>0\ \exists N\ \forall x\in A:\quad
\sum_{k=N}^\infty |x_k|^p<\eps.
\]
Deducir que $K=\{x\in\ell^2:0\le x_n\le1/n\}$ es compacto.
\end{ejercicio}


\begin{ejercicio}[Lectura 14, ejercicio 6]\label{ej:lectura-14-6}
Sea $W$ una vecindad compacta de $x$ y $f:W\to[0,\infty)$ continua, con $f(x)=0$ y $f(y)>0$ para $y\in W\setminus\{x\}$. Probar que para toda vecindad $V$ de $x$ existe $r>0$ tal que $\{y\in W:f(y)<r\}\subset V$.
\end{ejercicio}


En un compacto métrico siempre podemos extraer una subsucesión convergente. Para reconocer compactos de sucesiones o funciones conviene buscar dos ingredientes: control de finitas coordenadas y una estimación uniforme de la cola. El primero por sí solo no garantiza convergencia en la norma del espacio.

\section{Acotación total y caracterización métrica}\label{lectura:15}
En $\R^n$, ser cerrado y acotado equivale a ser compacto. En $\ell^\infty$ ya encontramos una bola cerrada que contiene infinitos puntos separados por una distancia fija. La noción que excluye esta situación es la acotación total: a cualquier escala, el espacio debe poder describirse con finitos centros. La completitud garantiza, además, que los límites que esa descripción produce permanezcan dentro del espacio. La caracterización secuencial obtenida en la sección \ref{lectura:14} nos permite precisar ahora esta relación entre aproximación finita y existencia de límites.

\begin{definicion}
Un espacio métrico $X$ es \emph{totalmente acotado}\index{acotacion total@acotación total} si para cada $\eps>0$ existen $x_1,\ldots,x_m\in X$ tales que $X=\bigcup_{i=1}^mB(x_i,\eps)$. Los centros constituyen una \emph{red finita de radio $\eps$}. Aquí ``red'' significa conjunto de centros, no una familia indexada por un conjunto dirigido.
\end{definicion}
Un espacio totalmente acotado no vacío es acotado: basta una cubierta por finitas bolas de radio $1$ y estimar las distancias de los centros a uno de ellos. Todo subconjunto $A$ de un espacio totalmente acotado es totalmente acotado con su métrica relativa. Para probarlo, se cubre el espacio con bolas de radio $\eps/2$ y, de cada bola que toca $A$, se elige un centro nuevo en $A$; las bolas de radio $\eps$ con los nuevos centros cubren $A$. Este cambio de centros evita suponer que los centros originales pertenecen a $A$.

\begin{teorema}\label{ch15:criterio}
Un espacio métrico es compacto si y solo si es completo y totalmente acotado.
\end{teorema}
\begin{proof}
Si es compacto, es completo por el ejercicio \ref{ej:lectura-13-3}. La cubierta por todas las bolas de un radio fijo tiene una subcubierta finita, lo que prueba la acotación total.

Para el recíproco, sea $(x_n)$ una sucesión. Una cubierta finita por bolas de radio $1/2$ contiene una bola a la que pertenecen infinitos términos: de lo contrario, la unión de las cantidades finitas de índices sería finita. Sea $I_1$ un conjunto infinito de esos índices. Cubrimos ahora $X$ por finitas bolas de radio $1/4$ y elegimos un conjunto infinito $I_2\subset I_1$ cuyos términos estén en una sola bola. Continuamos con radios $2^{-j}$. Elegimos índices $n_j\in I_j$ con $n_{j+1}>n_j$; esto es posible porque cada $I_j$ es infinito. Si $k,l\ge j$, entonces $n_k,n_l\in I_j$ y
\[
d(x_{n_k},x_{n_l})<2\cdot2^{-j}.
\]
La subsucesión es de Cauchy. La completitud le proporciona un límite en $X$. Así toda sucesión admite una subsucesión convergente y el teorema \ref{ch14:equivalencias} implica compacidad.
\end{proof}

\begin{observacion}[Heine--Borel y su alcance]
En $\R^n$, un conjunto acotado está contenido en una caja $[-M,M]^n$. Para cada $\eps>0$, dividimos cada intervalo en finitos subintervalos de longitud menor que $\eps/\sqrt n$. Cada caja pequeña tiene diámetro menor que $\eps$; eligiendo un punto del conjunto en cada caja que lo toque obtenemos una cubierta finita por bolas de radio $\eps$. Por tanto el conjunto es totalmente acotado. Si además es cerrado, es completo como subespacio cerrado del espacio completo $\R^n$, y el teorema da compacidad. La necesidad de cerrado y acotado ya se probó. En dimensión infinita, acotado no implica totalmente acotado, como muestran los vectores $e_n$ de la bola de $\ell^\infty$.\index{teorema de Heine--Borel}
\end{observacion}

\begin{proposicion}\label{ch15:relativa}
Para $A\subset X$ son equivalentes:
\begin{enumerate}[label=(\roman*)]
\item la cerradura $\cl A^{\,X}$ es compacta;
\item toda sucesión en $A$ tiene una subsucesión convergente a algún punto de $X$;
\item $A$ es totalmente acotado y $\cl A^{\,X}$ es completo.
\end{enumerate}
Cuando se cumplen, $A$ es \emph{relativamente compacto en $X$}.\index{compacidad!relativa}
\end{proposicion}
\begin{proof}
De (i) se sigue (ii) aplicando compacidad secuencial a la cerradura. También se sigue (iii), porque la cerradura es completa y totalmente acotada, y sus subconjuntos son totalmente acotados.

Para (ii)$\Rightarrow$(i), sea $(z_n)$ una sucesión en $\cl A$. Si $A$ es vacío no hay nada que probar. En otro caso, la definición de cerradura permite elegir $a_n\in A$ con $d(z_n,a_n)<1/n$. Por (ii), una subsucesión $a_{n_j}$ converge a $a\in X$. Entonces
\[
d(z_{n_j},a)\le 1/n_j+d(a_{n_j},a)\longrightarrow0.
\]
Como $\cl A$ es cerrado, $a\in\cl A$. Así $\cl A$ es secuencialmente compacto. Para (iii)$\Rightarrow$(i), el ejercicio \ref{ej:lectura-15-1} demuestra que $\cl A$ es totalmente acotado; se aplica el teorema \ref{ch15:criterio}.
\end{proof}

\begin{ejemplo}[El espacio ambiente importa]
El intervalo $(0,1)$ es relativamente compacto en $\R$, pues su cerradura allí es $[0,1]$. No es relativamente compacto en sí mismo: su cerradura relativa es $(0,1)$ y la sucesión $1/n$ no tiene ninguna subsucesión que converja a un punto de ese espacio. En un ambiente completo, ser relativamente compacto equivale a ser totalmente acotado. En un ambiente incompleto hay que comprobar también la completitud de la cerradura.
\end{ejemplo}

\section{Compacidad alrededor de cada punto}
\begin{definicion}
Un espacio métrico $X$ es \emph{localmente compacto}\index{compacidad!local} si cada punto tiene una vecindad compacta: para cada $x$ existen un abierto $O$ y un compacto $K$ con $x\in O\subset K$. Equivalentemente, cada punto tiene una vecindad cuya cerradura es compacta. Una vecindad no tiene por qué ser abierta.
\end{definicion}
La equivalencia se verifica así: una vecindad compacta es cerrada, por lo que es su propia cerradura; si una vecindad tiene cerradura compacta, esa cerradura sigue siendo una vecindad del punto.

\begin{proposicion}\label{ch15:subespacios}
Todo subespacio abierto o cerrado de un espacio métrico localmente compacto es localmente compacto.
\end{proposicion}
\begin{proof}
Sea $x\in A$ y sea $K$ una vecindad compacta de $x$ en $X$. Si $A$ es cerrado, $K\cap A$ es compacto y contiene una vecindad relativa de $x$ en $A$.

Si $A$ es abierto, elegimos un abierto $O$ con $x\in O\subset K$. Existe $r>0$ con $\overline B(x,r)\subset A\cap O$: basta escoger primero una bola de radio $2r$ contenida en ese abierto. La bola cerrada es un cerrado de $K$, luego es compacta, y constituye una vecindad de $x$ tanto en $X$ como en $A$.
\end{proof}

\begin{proposicion}\label{ch15:engrosamiento}
Si $X$ es localmente compacto y $A\subset X$ es compacto, existe $r>0$ tal que $\cl{B(A,r)}$ es compacto, donde $B(A,r)=\{x:\dist(x,A)<r\}$.
\end{proposicion}
\begin{proof}
Para $A=\varnothing$, interpretamos $B(A,r)=\varnothing$. En otro caso, por el ejercicio \ref{ej:lectura-15-3} cada $a\in A$ admite $r_a>0$ con $\overline B(a,r_a)$ compacta. Elegimos una subcubierta finita $B(a_i,r_i/2)$ de $A$, ponemos $r=\frac14\min_i r_i$ y $K=\bigcup_i\overline B(a_i,r_i)$. El conjunto $K$ es compacto. Si $z\in B(A,r)$, hay $a\in A$ con $d(z,a)<r$, y después un $i$ con $d(a,a_i)<r_i/2$. Así $d(z,a_i)<r+r_i/2<r_i$ y $z\in K$. Como $K$ es cerrado, también $\cl{B(A,r)}\subset K$. Esta cerradura es un cerrado de un compacto, luego compacta. No se ha supuesto que todas las bolas cerradas de $X$ sean compactas: solo se emplean radios locales adecuados.
\end{proof}

\section{Ejercicios propuestos: compacidad local y funciones de corte}
\begin{ejercicio}[Lectura 15, ejercicio 1]\label{ej:lectura-15-1}
Si $A\subset X$ es totalmente acotado, probar que su cerradura también lo es.
\end{ejercicio}


\begin{ejercicio}[Lectura 15, ejercicio 2]\label{ej:lectura-15-2}
Probar que el recorrido de toda sucesión de Cauchy es totalmente acotado.
\end{ejercicio}


\begin{ejercicio}[Lectura 15, ejercicio 3]\label{ej:lectura-15-3}
Probar que $X$ es localmente compacto si y solo si cada $x$ admite $r_x>0$ con la bola cerrada $\overline B(x,r_x)$ compacta.
\end{ejercicio}


\begin{ejercicio}[Lectura 15, ejercicio 4]\label{ej:lectura-15-4}
Probar que $\Q$, con la métrica inducida de $\R$, no es localmente compacto.
\end{ejercicio}


\begin{ejercicio}[Lectura 15, ejercicio 5]\label{ej:lectura-15-5}
Sean $X$ localmente compacto, $A\subset X$ compacto y $U\subset X$ abierto con $A\subset U$. Probar:
\begin{enumerate}[label=(\roman*)]
\item existen $V$ abierto y $B$ compacto con $A\subset V\subset B\subset U$;
\item existe $f:X\to[0,1]$ continua, igual a $1$ en $A$, cuyo soporte es compacto y está contenido en $U$.
\end{enumerate}
\end{ejercicio}


\begin{ejercicio}[Lectura 15, ejercicio 6]\label{ej:lectura-15-6}
Si el compacto $A$ está cubierto por abiertos $U_1,\ldots,U_n$ de un espacio localmente compacto $X$, construir funciones continuas $f_k:X\to[0,1]$ de soporte compacto contenido en $U_k$, tales que $\sum_k f_k\le1$ en $X$ y $\sum_kf_k=1$ en $A$.
\end{ejercicio}


\section{De la extracción finita a la localización}
Las cubiertas, las subsucesiones y las redes finitas son tres formas complementarias de reconocer la compacidad métrica. Las primeras convierten datos locales en controles uniformes; las segundas permiten obtener límites de aproximaciones. La acotación total controla todas las escalas; la completitud retiene sus límites. La compacidad local permite construir funciones de corte y particiones finitas de la unidad, herramientas que más adelante harán posible localizar problemas de integración y aproximación.
